\section{模型的评价与推广}
\subsection{模型的优点和不足}
模型的主要优点可归纳为以下三点。
\begin{enumerate}
 \item \textbf{运行与账单逐段对应。}能量平衡、损耗和库存约束统一建模，真实反馈与交易结算独立核验；指定日期和全年费用均可追溯至十分钟记录。
 \item \textbf{信息条件清楚。}预测只用已完成观测和已发布预报，跨日库存连续承接；信息子集对照区分了新预报与重规划机会的作用。
 \item \textbf{求解与解释较简洁。}储能规划采用线性模型，日内规模随剩余时域缩短；设备、分位和合约检验分别解释费用变化，问题一另有整数模型复核。
\end{enumerate}

模型的不足主要体现在以下三个方面。
\begin{enumerate}
 \item \textbf{风险与库存价值处理较简化。}逐段残差分位未联合刻画时间相关性，也不等于可靠率；日末恢复目标与当前缺口优先反馈尚未联合优化跨日库存价值。
 \item \textbf{参数具有样本依赖。}1月短窗口选出的分位未必全年最省，逐月重选与所测试的滚动反馈也未改善累计费用。现有经济结果属于2025年再分析，不能直接外推。
 \item \textbf{实际运行条件尚未完整纳入。}模型未计电池折损、外网容量及响应延迟，并假定紧急电可补足缺口；追加购电是否允许撤回，也会改变费用与策略。
\end{enumerate}

\subsection{模型的推广}
后续推广可沿数据、设备和控制三个方向逐步开展。
\begin{enumerate}
 \item \textbf{其他年度与负荷类型。}替换负荷、光伏和电价输入，按实际发布时间对齐；在未参与开发的数据上，冻结模型族与选参规则后连续回测，同时比较费用、紧急电量和边界库存。
 \item \textbf{供电限制与寿命管理。}有设备资料后，加入购电功率上限、循环吞吐成本和随健康状态变化的容量。紧急电也受限时需记录缺供量并设置代价；扩容评价还应计入投资成本。
 \item \textbf{多时段风险控制。}构造保留误差时间相关性的场景或不确定集合，在跨日滚动中估计末端库存价值。相关研究\upcite{li2024mpc,li2025dro}可提供方法背景，仍须在相同信息、合约和初始库存下比较经济效果与计算开销。这些方向属于待验证扩展。
\end{enumerate}
\label{end:body}
